\chapter{第六讲\quad 客观要件二：结果}
\label{lecture:06}
\noindent{\small\hyperref[ref:lecture:06]{查阅本讲扩展规则}}\par

本讲把三个不同问题拆开。罪状要求何种危险或实害。基本犯达到什么阶段。加重结果能否进入该基本犯的法定刑升格条件。

\section{第一节\quad 法益侵害事实的分类}

\paragraph{定义与法条入口}
对罪名 $j$，令 $B_j$ 表示罪状要求的行为及其其他基础条件，$D_j$ 表示个案达到法定危险程度，$H_j$ 表示实害达到法定门槛，$Q_j$ 表示已满足所选罪状的客观结果或危险条件。把罪状依是否要求个案证明危险或实害分三型：
\[
\begin{array}{c|c|l}
\text{教材分类} & \text{客观核心条件} & \text{对应条文示例}\\ \hline
\text{抽象危险型} & B_j & \text{刑法第141条基本档}\\
\text{具体危险型} & B_j\land D_j & \text{刑法第114条}\\
\text{实害型} & B_j\land H_j & \text{刑法第142条}
\end{array}
\]
教材使用「行为犯」强调基本构成不以实害为门槛，使用「结果犯」强调法定实害条件。具体危险型虽然不要求实害，仍须证明个案危险。分类应具体到条文和处罚档，而非把一个罪名的所有档次一概而论。第141条现行罪名为生产、销售、提供假药罪，提供分支另有人员与明知要求。

\paragraph{推导}
甲明知以面粉冒充药品并生产、销售，题设给定物品依法属于假药，相关主体、行为与主观条件已证明，没有严重实害，也没有需要证明的个案危险。
\[
\begin{aligned}
\Gamma&\vdash B_{141}\land\neg H_{141},\\
B_{141}&\Longrightarrow Q_{141,\mathrm{base}},\\
\Gamma&\vdash Q_{141,\mathrm{base}} .
\end{aligned}
\]
因此无严重实害不妨碍满足第141条基本档的客观要求。若物品依法是劣药而非假药，须重新套用第142条。若已证明不曾造成法条所需严重危害，则
$B_{142}\land\neg H_{142}\Rightarrow \neg(B_{142}\land H_{142})$。不能据此认定第142条的实害既遂。其他罪名或形态应各审其要件。

\paragraph{边界}
燃烧了物品只说明有燃烧行为。若火在完全隔绝环境中，证据已排除危害公共安全的具体危险，则 $D_{114}$ 不成立。若改为火势已经威胁多人聚集的建筑、后来及时被扑灭，则可能证明 $D_{114}$，无需等到有人实际重伤。是否达到危险门槛应说明人员暴露、传播条件和危险程度，不能把「事后未出事」当作安全证明。

\section{第二节\quad 结果加重犯}

\paragraph{构造}
令 $A$ 为基本犯，$a$ 为其相关实行行为，$r$ 为待加重的结果。设 $B_A(a)$ 为基本犯罪的相关前提，$L_A(r)$ 为适用的刑法条文或有权解释规定该结果进入加重档，$R(r)$ 为实际结果，$G_A(a,r)$ 为该结果由基本实行行为所创设的风险实现且可归责，$F_A(r)$ 为达到所要求的结果罪过。教材模型为
\[
\mathrm{Aggr}_A(a,r)
  = B_A(a)\land L_A(r)\land R(r)\land G_A(a,r)\land F_A(r).
\]
输出是同一基本罪名下的加重档。$F_A$ 的内容随条文、手段与所采解释而定，不能一律写成过失。基本行为、后续独立行为及其目的，应按时间分开。即使出现加重结果，若 $\neg L_A(r)$，也要回到普通竞合规则，核对其他罪的要件后才可能择一重处或并罚。

\paragraph{推导}
甲持刀刺伤乙腿部，故意在于伤害，乙因刺伤失血死亡。题设给定甲对死亡至少有可归责的过失，鉴定确认伤口出血连续导致死亡。按刑法第二百三十四条第二款，
\[
\begin{aligned}
\text{伤害故意与刺伤}&\Longrightarrow B_{\text{伤害}}(a),\\
\text{死亡与连续失血机制}&\Longrightarrow R(r)\land G_{\text{伤害}}(a,r),\\
\text{法条与已证过失}&\Longrightarrow L_{\text{伤害}}(r)\land F_{\text{伤害}}(r).
\end{aligned}
\]
合取得出故意伤害致人死亡的加重档。若甲对死亡具有杀人故意，则应重新审查故意杀人规范。若乙已稳定、后在途中因独立事故死亡，且该事故单独决定死亡，则 $G_{\text{伤害}}(a,r)$ 可能不成立。时间上在伤害之后发生，不能替代机制判断。

\paragraph{目的与条文的双重核对}
抢劫中为压制反抗而实施的暴力，可能是第二百六十三条第五项加重结果的基本实行行为。取财已结束，甲另因私怨伤人，设这一行为为 $b$。即使 $\mathrm{Cause}(b,r)$，仍须证明 $b$ 是抢劫的实行行为才可将 $r$ 归入抢劫加重档。在题设明确私怨暴力独立的情形，应分别评价两项行为。

教材常用的法条检查包括第二百三十八条非法拘禁致重伤、死亡，并与使用暴力致伤残、死亡的转致规定区别。另查第二百三十六条强奸、第二百四十条拐卖、第二百五十七条暴力干涉婚姻自由、第二百六十条虐待及第一百一十五条危害公共安全的结果档。抢夺案件还应核对两高关于办理抢夺刑事案件的司法解释，区别重伤、死亡情节与特定驾驶夺取情形按抢劫定罪的规则。

\paragraph{边界}
刑法修正案（九）改变了第二百三十九条的加重表述。现行款列杀害被绑架人，或故意伤害被绑架人致重伤、死亡。单凭捆绑时不慎造成窒息死亡，不能推出满足这一款的特别条件。若绑架和过失致死的各项要件均成立，且系同一行为，教材再按想象竞合分析。还要核对行为时法及从旧兼从轻。另如遗弃后发生死亡，也不能为第二百六十一条自行补造一个死亡加重档。

\section*{闭卷记忆与自测}

\paragraph{复述路线}
看到分则条件先写「行为、个案危险、实害」，圈出条文实际要求的证明对象。遇到升格，背「基、条、果、链、过」，依次还原为 $B_A,L_A,R,G_A,F_A$。最后问一句「致害的是哪一个实行行为」，并写出该行为时的目的。

\paragraph{题一}
销售劣药，题设确认未对健康造成严重危害。能否用「药品具有抽象危险」补足第142条？

\paragraph{核对}
不能。第142条要求法定的严重危害，题设令 $H_{142}=0$。第141条基本档的证明对象不能移植到第142条。

\paragraph{题二}
甲抢劫得手并结束后，又独立实施报复伤害致人重伤。第二百六十三条虽有致人重伤加重规定，为什么仍需另查伤害行为？

\paragraph{核对}
$L_A(r)$ 只说明条文承认这种结果类型。此案重伤来自独立报复行为，不具题设所需的抢劫实行行为和结果归责链。应分别评价抢劫与伤害。
